\section*{Problems}

\subsection*{Problem statements}

% Remember
\begin{problema}
	\label{prob:06f83cc}
	A telephone call center receives calls independently at a constant average rate $\lambda$ per minute. State, without performing any calculation, the probability mass function $P(X=k)$ of $X\sim\text{Pois}(\lambda)$ and its support $\mathcal{S}_X$.
\end{problema}

% Understand
\begin{problema}
	\label{prob:c35cfa7}
	An industrial system records machine failures at a rate of $0.5$ failures/hour, under a Poisson process. Explain why the number of failures in a $4$-hour shift follows $\text{Pois}(4\times0.5)=\text{Pois}(2)$, and why modeling the same shift as $n=4$ independent hourly Bernoulli trials with $p=1-e^{-0.5}$ produces exactly the same probability of at least one failure, $P(X\ge1)=1-e^{-2}$.
\end{problema}

% Apply
\begin{problema}
	\label{prob:8fd3390}
	The number of people arriving per day at an emergency room follows $\text{Pois}(\lambda=5)$.
	\begin{enumerate}
		\item Compute $P(X\le3)$, the probability that at most $3$ people arrive in one day.
		\item Compute $P(X\ge8)$, the probability that at least $8$ people arrive.
	\end{enumerate}
\end{problema}

% Analyze
\begin{problema}
	\label{prob:8dad711}
	Let $X\sim\text{Pois}(\lambda_1)$ and $Y\sim\text{Pois}(\lambda_2)$ be independent (you may use without proof that $X+Y\sim\text{Pois}(\lambda_1+\lambda_2)$). Prove that the conditional distribution of $X$ given $X+Y=n$ is Binomial with parameters $n$ and $p=\lambda_1/(\lambda_1+\lambda_2)$, and apply the result with $\lambda_1=3$, $\lambda_2=5$, $n=8$ to compute $P(X=3\mid X+Y=8)$.
\end{problema}

% Evaluate
\begin{problema}
	\label{prob:5e9408a}
	Let $X_1,\dots,X_{100}$ be i.i.d. $\text{Pois}(\lambda=2)$, with $S_{100}=\sum X_i \sim \text{Pois}(200)$. By the Central Limit Theorem, $S_{100}$ is approximated by $N(200,200)$. Evaluate the quality of this approximation by computing $P(S_{100}\le220)$ exactly (through the summed Poisson distribution) and through the Normal approximation with continuity correction, and determine whether the absolute error is small enough to justify practical use of the approximation.
\end{problema}

% Create
\begin{problema}
	\label{prob:9dc367e}
	Design an original example (context and rate $\lambda$) that is naturally modeled by a Poisson distribution, different from the call-center, manufacturing, or web-server examples already seen in this section. State the probability question of interest and compute $\mathbb{E}[X]$ and $\text{Var}(X)$ for the parameter you chose.
\end{problema}

\subsection*{Detailed solutions}

% Remember
\begin{solproblema}[prob:06f83cc]
	$P(X=k) = \dfrac{e^{-\lambda}\lambda^k}{k!}$, for $k \in \mathcal{S}_X = \set{0,1,2,\dots}$.
\end{solproblema}

% Understand
\begin{solproblema}[prob:c35cfa7]
	In a Poisson process, the number of events in a time interval scales linearly with its duration: if the rate is $0.5$ failures/hour, the expected rate over $4$ hours is $4\times0.5=2$, so the count follows $\text{Pois}(2)$, and $P(X\ge1)=1-e^{-2}\approx0.8647$. When discretizing the same shift into $4$ independent hourly Bernoulli trials, the probability of at least one failure in a single hour is $p=1-e^{-0.5}\approx0.3935$ (the complement of $P(\text{Pois}(0.5)=0)$), and the probability of at least one failure in the $4$ hours is $1-(1-p)^4 = 1-(e^{-0.5})^4 = 1-e^{-2}$, identical to the previous result. Both models coincide because they are two equivalent ways to count the same rare-arrivals process: one in continuous time (Poisson) and the other discretized into subintervals (Bernoulli/Binomial), and both share the same exact probability of "zero events" in each subinterval.
\end{solproblema}

% Apply
\begin{solproblema}[prob:8fd3390]
	With $\lambda=5$:
	\begin{enumerate}
		\item $P(X\le3) = e^{-5}\left(1+5+\dfrac{25}{2}+\dfrac{125}{6}\right) \approx 0.2650$.
		\item $P(X\ge8) = 1-P(X\le7) \approx 1-0.8670 = 0.1330$.
	\end{enumerate}
\end{solproblema}

% Analyze
\begin{solproblema}[prob:8dad711]
	By the definition of conditional probability and independence:
	\begin{align*}
		P(X=k\mid X+Y=n) = \dfrac{P(X=k)P(Y=n-k)}{P(X+Y=n)} = \dfrac{\dfrac{e^{-\lambda_1}\lambda_1^k}{k!}\cdot\dfrac{e^{-\lambda_2}\lambda_2^{n-k}}{(n-k)!}}{\dfrac{e^{-(\lambda_1+\lambda_2)}(\lambda_1+\lambda_2)^n}{n!}} = \comb{n}{k}\left(\dfrac{\lambda_1}{\lambda_1+\lambda_2}\right)^k\left(\dfrac{\lambda_2}{\lambda_1+\lambda_2}\right)^{n-k},
	\end{align*}
	which is the PMF of $\text{Bin}(n,p)$ with $p=\lambda_1/(\lambda_1+\lambda_2)$. With $\lambda_1=3,\lambda_2=5,p=3/8$: $P(X=3\mid X+Y=8) = \comb{8}{3}\left(\dfrac{3}{8}\right)^3\left(\dfrac{5}{8}\right)^5 \approx 0.2815$.
\end{solproblema}

% Evaluate
\begin{solproblema}[prob:5e9408a]
	Exact: $P(S_{100}\le220) = \sum_{k=0}^{220}\dfrac{e^{-200}200^k}{k!} \approx 0.9862$ (reference value under $\text{Pois}(200)$). Normal approximation with continuity correction: $P(S_{100}\le220) \approx \Phi\!\left(\dfrac{220.5-200}{\sqrt{200}}\right) = \Phi(1.4496) \approx 0.9265$. The absolute error between the two methods is considerable ($\approx0.06$), larger than one might expect for $\lambda=200$; this indicates that, although $n=100$ is large, the Normal approximation in the tail (where the exact probability is already high) may be less accurate than near the center of the distribution, and it is preferable to use the exact Poisson calculation (or statistical software) instead of the Normal approximation when tail precision is required.
\end{solproblema}

% Create
\begin{solproblema}[prob:9dc367e]
	\textbf{Example:} a library records an average of $\lambda=4$ interlibrary-loan requests per day, modeled as a Poisson process. Question of interest: what is the probability that more than $6$ requests are received on a given day, $P(X>6)$? The moments are $\mathbb{E}[X]=\text{Var}(X)=\lambda=4$. The probability would be computed as $P(X>6) = 1-\sum_{k=0}^{6}\dfrac{e^{-4}4^k}{k!}$.
\end{solproblema}
